% GENERATED FILE. Edit canonical lessons, then run npm run generate:books.
% canonical-source-hash: fnv1a64:fdd55b5a5d8f2732
% canonical-lessons: JA-C259-yori, JA-C259-dochira, JA-C259-ichiban, JA-C259-nijikan-gurai, JA-R259-comparing-and-how-long

\chapter{Comparing, and How Long}
\label{ch:ja-comparing-and-how-long}
\hlchaptermodality{259}

\begin{chapteropening}
\textbf{By the end of this chapter:} \emph{I can compare two things (hikōki wa densha yori hayai desu; dochira ga suki desu ka), say which is the most (ichiban), and say how long something takes and from when to when (nijikan gurai; kuji kara goji made).}

Complete JA-R259-comparing-and-how-long: compare two places, say which you like best and how long the trip takes.
\end{chapteropening}

\section[\texorpdfstring{\ja{より} (yori)}{より (yori)}]{\ja{より} (yori) — than}

\label{lesson:JA-C259-yori}



\begin{quote}
Say \textquotedblleft{}fast\textquotedblright{} and \textquotedblleft{}a plane\textquotedblright{}. (\emph{hayai}, \emph{hikōki}.) Then \textquotedblleft{}you had better rest\textquotedblright{}. (\emph{yasunda hō ga ii desu}.) Its \textbf{\ja{ほう}} is back today.
\end{quote}

\subsection*{You'll want to know: A \ja{は} B \ja{より}}
\begin{quote}
\textbf{\ja{ひこうきは} \ja{でんしゃより} \ja{はやいです。}}
\end{quote}

\begin{quote}
\emph{hikōki wa densha yori hayai desu.} — \textquotedblleft{}A plane is faster than a train.\textquotedblright{}
\end{quote}

\textbf{\ja{より}} (\emph{yori}), \textquotedblleft{}than\textquotedblright{}, follows the thing you measure against. The adjective does not change: \textbf{\ja{はやい}} is \textquotedblleft{}faster\textquotedblright{} here.

\begin{quote}
\textbf{\ja{なつは} \ja{ふゆより} \ja{あついです。}}
\end{quote}

\begin{quote}
\emph{natsu wa fuyu yori atsui desu.} — \textquotedblleft{}Summer is hotter than winter.\textquotedblright{}
\end{quote}

\subsection*{You'll want to know: B \ja{より} A \ja{の} \ja{ほうが}}
To put the winner last, \textbf{\ja{の} \ja{ほうが}} (\emph{no hō ga}), \textquotedblleft{}the side of\textquotedblright{}, marks it:

\begin{quote}
\textbf{\ja{ふゆより} \ja{なつの} \ja{ほうが} \ja{すきです。}}
\end{quote}

\begin{quote}
\emph{fuyu yori natsu no hō ga suki desu.} — \textquotedblleft{}I like summer more than winter.\textquotedblright{}
\end{quote}

\begin{quote}
\textbf{\ja{やまより} \ja{うみの} \ja{ほうが} \ja{すきです。}}
\end{quote}

\begin{quote}
\emph{yama yori umi no hō ga suki desu.} — \textquotedblleft{}I like the sea more than the mountains.\textquotedblright{}
\end{quote}

\subsection*{Your turn}
Say these aloud:

\begin{itemize}
\raggedright
  \item \textquotedblleft{}I like the sea more than the mountains\textquotedblright{} — \emph{yama yori umi no hō ga suki desu}
  \item \textquotedblleft{}I like the river more than the mountains\textquotedblright{} — \emph{yama yori kawa no hō ga suki desu}
  \item \textquotedblleft{}I like the town more than the mountains\textquotedblright{} — \emph{yama yori machi no hō ga suki desu}
\end{itemize}

\begin{tcolorbox}[breakable,colback=teal!4,colframe=teal!35!black,title={Before you move on}]
\textquotedblleft{}Than\textquotedblright{}? (\textbf{\ja{より}}.) Does \textbf{\ja{はやい}} change for \textquotedblleft{}faster\textquotedblright{}? (No.) \textquotedblleft{}I like summer more than winter\textquotedblright{}? (\textbf{\ja{ふゆより} \ja{なつの} \ja{ほうが} \ja{すきです}}.)
\end{tcolorbox}
\section[\texorpdfstring{\ja{どちら} (dochira)}{どちら (dochira)}]{\ja{どちら} (dochira) — which (of two)}

\label{lesson:JA-C259-dochira}



\begin{quote}
Before the new word: say \textquotedblleft{}I like the sea more than the mountains\textquotedblright{}. (\emph{yama yori umi no hō ga suki desu}.) Then \textquotedblleft{}where\textquotedblright{} and \textquotedblleft{}which\textquotedblright{}. (\emph{doko}, \emph{dono}.)
\end{quote}

\subsection*{You'll want to know: \ja{どちら}}
\textbf{\ja{どちら}} — \emph{dochira} — \textquotedblleft{}which\textquotedblright{} of two. Name the two with \textbf{\ja{と}}, then ask:

\begin{quote}
\textbf{\ja{なつと} \ja{ふゆと} \ja{どちらが} \ja{すきですか。}}
\end{quote}

\begin{quote}
\emph{natsu to fuyu to dochira ga suki desu ka.} — \textquotedblleft{}Which do you like better, summer or winter?\textquotedblright{}
\end{quote}

\begin{quote}
\textbf{\ja{なつの} \ja{ほうが} \ja{すきです。}}
\end{quote}

\begin{quote}
\emph{natsu no hō ga suki desu.} — \textquotedblleft{}I like summer better.\textquotedblright{}
\end{quote}

\begin{cousinweb}
\textbf{\ja{どちら}} (\emph{dochira}) is a native Japanese word (wago), inherited rather than borrowed from Chinese: the \emph{do-} of \textbf{\ja{どこ}} and \textbf{\ja{どの}}, \textquotedblleft{}which\textquotedblright{}, with \emph{-chira}, a direction ending. JMdict marks it as usually written in kana and glosses it \textquotedblleft{}which one (esp. of two alternatives)\textquotedblright{} and \textquotedblleft{}which way\textquotedblright{} (JMdict).
\end{cousinweb}

\subsection*{Your turn}
Say these aloud:

\begin{itemize}
\raggedright
  \item \emph{dochira}
  \item \textquotedblleft{}Which do you like better, the sea or the mountains?\textquotedblright{} — \emph{umi to yama to dochira ga suki desu ka}
  \item \textquotedblleft{}Which do you like better, tea or coffee?\textquotedblright{} — \emph{ocha to kōhī to dochira ga suki desu ka}
\end{itemize}

\begin{tcolorbox}[breakable,colback=teal!4,colframe=teal!35!black,title={Before you move on}]
\textquotedblleft{}Which\textquotedblright{}, of two? (\textbf{\ja{どちら}}, \emph{dochira}.) And the answer starts? (With the winner and \textbf{\ja{の} \ja{ほうが}}.)
\end{tcolorbox}
\section[\texorpdfstring{\ja{いちばん} (ichiban)}{いちばん (ichiban)}]{\ja{いちばん} (ichiban) — the most; number one}

\label{lesson:JA-C259-ichiban}



\begin{quote}
Before the new word: ask \textquotedblleft{}which do you like better, summer or winter?\textquotedblright{} (\emph{natsu to fuyu to dochira ga suki desu ka}.) Then say \textquotedblleft{}summer is hotter than winter\textquotedblright{}. (\emph{natsu wa fuyu yori atsui desu}.) And \textquotedblleft{}one\textquotedblright{}. (\emph{ichi}.)
\end{quote}

\subsection*{You'll want to know: \ja{いちばん}}
\textbf{\ja{いちばん}} — \emph{ichiban} — \textquotedblleft{}the most\textquotedblright{}, \textquotedblleft{}number one\textquotedblright{}. It goes before the adjective, and with more than two things to choose from:

\begin{quote}
\textbf{\ja{なつが} \ja{いちばん} \ja{すきです。}}
\end{quote}

\begin{quote}
\emph{natsu ga ichiban suki desu.} — \textquotedblleft{}I like summer best.\textquotedblright{}
\end{quote}

\begin{quote}
\textbf{\ja{かぞくの} \ja{なかで} \ja{ちちが} \ja{いちばん} \ja{たかいです。}}
\end{quote}

\begin{quote}
\emph{kazoku no naka de chichi ga ichiban takai desu.} — \textquotedblleft{}In my family my father is the tallest.\textquotedblright{}
\end{quote}

\begin{cousinweb}
\textbf{\ja{いちばん}} (\emph{ichiban}) is Sino-Japanese (kango): two Chinese characters read with their on'yomi, readings borrowed from Chinese. They are \emph{ichi}, \textquotedblleft{}one\textquotedblright{} (Middle Chinese \emph{'jit}) and \emph{ban}, \textquotedblleft{}number in a series\textquotedblright{} (\emph{phjon}), so \textquotedblleft{}number one\textquotedblright{} (KANJIDIC2; Middle Chinese in Baxter's transcription). JMdict gives \textquotedblleft{}number one; first\textquotedblright{} and \textquotedblleft{}best; most\textquotedblright{} (JMdict).
\end{cousinweb}

\subsection*{Your turn}
Say these aloud:

\begin{itemize}
\raggedright
  \item \emph{ichiban}
  \item \textquotedblleft{}I like spring best\textquotedblright{} — \emph{haru ga ichiban suki desu}
  \item \textquotedblleft{}I like autumn best\textquotedblright{} — \emph{aki ga ichiban suki desu}
\end{itemize}

\begin{tcolorbox}[breakable,colback=teal!4,colframe=teal!35!black,title={Before you move on}]
\textquotedblleft{}The most\textquotedblright{}? (\textbf{\ja{いちばん}}, \emph{ichiban}.) Two things to compare: \textbf{\ja{どちら}} or \textbf{\ja{いちばん}}? (\textbf{\ja{どちら}}; \textbf{\ja{いちばん}} is for more than two.)
\end{tcolorbox}
\section[\texorpdfstring{\ja{にじかんぐらい} … (nijikan gurai …)}{にじかんぐらい … (nijikan gurai …)}]{\ja{にじかんぐらい} / \ja{くじから} \ja{ごじまで} (nijikan gurai / kuji kara goji made) — how long}

\label{lesson:JA-C259-nijikan-gurai}



\begin{quote}
Say \textquotedblleft{}time, an hour\textquotedblright{} and \textquotedblleft{}a minute\textquotedblright{}. (\emph{jikan}, \emph{fun}.) Then \textquotedblleft{}to take\textquotedblright{}, of time. (\emph{kakaru}.)
\end{quote}

\subsection*{You'll want to know: \ja{〜}\ja{じかんぐらい}}
\begin{quote}
\textbf{\ja{くうこうまで} \ja{にじかん} \ja{かかります。}}
\end{quote}

\begin{quote}
\emph{kūkō made nijikan kakarimasu.} — \textquotedblleft{}It takes two hours to the airport.\textquotedblright{}
\end{quote}

A number before \textbf{\ja{じかん}} counts hours: \textbf{\ja{にじかん}} (\emph{nijikan}), two hours. Before \textbf{\ja{ふん}} it counts minutes: \textbf{\ja{ごふん}} (\emph{gofun}), five minutes. \textbf{\ja{ぐらい}} (\emph{gurai}), \textquotedblleft{}about\textquotedblright{}, follows the time:

\begin{quote}
\textbf{\ja{くうこうまで} \ja{にじかんぐらい} \ja{かかります。}}
\end{quote}

\begin{quote}
\emph{kūkō made nijikan gurai kakarimasu.} — \textquotedblleft{}It takes about two hours to the airport.\textquotedblright{}
\end{quote}

\begin{quote}
\textbf{\ja{えきまで} \ja{ごふんぐらい} \ja{あるきます。}}
\end{quote}

\begin{quote}
\emph{eki made gofun gurai arukimasu.} — \textquotedblleft{}I walk about five minutes to the station.\textquotedblright{}
\end{quote}

\subsection*{You'll want to know: \ja{〜}\ja{から} \ja{〜}\ja{まで}}
\textbf{\ja{から}} (\emph{kara}) after a time or a place is \textquotedblleft{}from\textquotedblright{}, and \textbf{\ja{まで}} is \textquotedblleft{}until\textquotedblright{}, \textquotedblleft{}as far as\textquotedblright{}:

\begin{quote}
\textbf{\ja{くじから} \ja{ごじまで} \ja{はたらきます。}}
\end{quote}

\begin{quote}
\emph{kuji kara goji made hatarakimasu.} — \textquotedblleft{}I work from nine to five.\textquotedblright{}
\end{quote}

\begin{quote}
\textbf{\ja{げつようびから} \ja{きんようびまで} \ja{はたらきます。}}
\end{quote}

\begin{quote}
\emph{getsuyōbi kara kin'yōbi made hatarakimasu.} — \textquotedblleft{}I work from Monday to Friday.\textquotedblright{}
\end{quote}

\subsection*{Your turn}
Say these aloud:

\begin{itemize}
\raggedright
  \item \textquotedblleft{}I work from nine to five\textquotedblright{} — \emph{kuji kara goji made hatarakimasu}
  \item \textquotedblleft{}I work from eight to five\textquotedblright{} — \emph{hachiji kara goji made hatarakimasu}
  \item \textquotedblleft{}I work from ten to five\textquotedblright{} — \emph{jūji kara goji made hatarakimasu}
\end{itemize}

\begin{tcolorbox}[breakable,colback=teal!4,colframe=teal!35!black,title={Before you move on}]
\textquotedblleft{}About two hours\textquotedblright{}? (\textbf{\ja{にじかんぐらい}}.) \textquotedblleft{}Five minutes\textquotedblright{}? (\textbf{\ja{ごふん}}.) \textquotedblleft{}From Monday to Friday\textquotedblright{}? (\textbf{\ja{げつようびから} \ja{きんようびまで}}.)
\end{tcolorbox}
\section[(review)]{Review — comparing, and how long}

\label{lesson:JA-R259-comparing-and-how-long}



\begin{quote}
Say \textquotedblleft{}a plane is faster than a train\textquotedblright{}. (\emph{hikōki wa densha yori hayai desu}.) Then \textquotedblleft{}about two hours\textquotedblright{}. (\emph{nijikan gurai}.)
\end{quote}

\subsection*{Your turn}
Say these aloud:

\begin{itemize}
\raggedright
  \item Summer is hotter than winter.
  \item I like the sea more than the mountains.
  \item In my family my father is the tallest.
  \item I walk about five minutes to the station.
  \item I work from Monday to Friday.
  \item I have been to the sea, and in the future I want to become a doctor.
\end{itemize}

\begin{tcolorbox}[breakable,colback=teal!4,colframe=teal!35!black,title={Before you move on}]
\textquotedblleft{}Than\textquotedblright{}? (\textbf{\ja{より}}.) \textquotedblleft{}Which, of two\textquotedblright{}? (\textbf{\ja{どちら}}.) \textquotedblleft{}The most\textquotedblright{}? (\textbf{\ja{いちばん}}.) \textquotedblleft{}About\textquotedblright{}? (\textbf{\ja{ぐらい}}.) \textquotedblleft{}From … until …\textquotedblright{}? (\textbf{\ja{から} … \ja{まで}}.)
\end{tcolorbox}
